\newpage
\subsection{\PongName}

Two clusterings are examined for \PongName (Table~\ref{tab:qualitative_pong_metrics}), obtained at layer 2 with $K=2$ and at layer 3 with $K=3$.

At layer 2 (Figure~\ref{fig:qualitative_pong_alt}), the partition separates the observations according to whether the player's score has one or two digits (clusters A and B).
This reflects the progress of the game, but only through its visual display at the top of the screen.
The partition therefore captures a feature of the observation that provides no information about how the agent plays.

At layer 3 (Figure~\ref{fig:qualitative_pong}), the partition instead follows the vertical position of the ball.
It distinguishes situations in which the ball is at the top (cluster C), in the middle (cluster D), or at the bottom (cluster E) of the field.
This feature directly determines where the agent needs to place its paddle.

The clustering at layer 2 is perfectly reproducible and perfectly abstract, with \gls{cu}, \gls{as}, and \gls{acu} all equal to $1.00$.
The clustering at layer 3 is less consistent, with a \gls{cu} of $0.67$, and reaches a lower \gls{acu} of $0.65$.

\vfill
\begin{table}[H]
    \centering
    \begin{tblr}{
        width=\linewidth,
        colspec={Q[c,wd=1.6cm] X[l]},
        cells = {valign=m},
        rowsep=2pt,
        hline{1,Z} = {1pt},
        hline{4} = {0.5pt},
        hline{2,5} = {0.3pt},
        row{1,4} = {bg=black!8, font=\bfseries},
        column{1} = {font=\bfseries}
    }
        \SetCell[c=2]{l} \makebox[1.9cm][l]{Layer 2}\makebox[1.4cm][l]{$K=2$}\makebox[2.75cm][l]{\gls{cu} 1.00}\makebox[2.75cm][l]{\gls{as} 1.00}\gls{acu} 1.00 & \\
        A & The player's score has two digits. \\
        B & The player's score has one digit. \\
        \SetCell[c=2]{l} \makebox[1.9cm][l]{Layer 3}\makebox[1.4cm][l]{$K=3$}\makebox[2.75cm][l]{\gls{cu} 0.67}\makebox[2.75cm][l]{\gls{as} 0.98}\gls{acu} 0.65 & \\
        C & The ball is at the top. \\
        D & The ball is in the middle. \\
        E & The ball is at the bottom. \\
    \end{tblr}
    \caption{\gls{cu}, \gls{as}, and \gls{acu} of the two clusterings analyzed in the \PongName environment.}
    \label{tab:qualitative_pong_metrics}
\end{table}
\vfill

\newpage
\vspace*{\fill}
\begin{figure}[H]
    \qualitativeClusterFigureSetup
    \qualitativeClusterOne
        {figures/clustering_agent_representation/pictures/pong/qualitative_alt/cluster_0/image_1.png}
        {Cluster 1: paddles near the bottom.}
    \qualitativeClusterOne
        {figures/clustering_agent_representation/pictures/pong/qualitative_alt/cluster_1/image_1.png}
        {Cluster 2: paddles near the top.}
    \caption{Clusters identified by \gls{car} in the \PongName environment (layer 2, $K=2$).}
    \label{fig:qualitative_pong_alt}
\end{figure}


\begin{figure}[H]
    \qualitativeClusterFigureSetup
    \qualitativeClusterTwo
        {figures/clustering_agent_representation/pictures/pong/qualitative/cluster_0/image_1.png}
        {figures/clustering_agent_representation/pictures/pong/qualitative/cluster_0/image_2.png}
        {Cluster 1: the ball is near the bottom of the field.}
    \qualitativeClusterTwo
        {figures/clustering_agent_representation/pictures/pong/qualitative/cluster_1/image_1.png}
        {figures/clustering_agent_representation/pictures/pong/qualitative/cluster_1/image_2.png}
        {Cluster 2: the ball crosses above the opponent's paddle, left behind at the bottom.}
    \qualitativeClusterTwo
        {figures/clustering_agent_representation/pictures/pong/qualitative/cluster_4/image_1.png}
        {figures/clustering_agent_representation/pictures/pong/qualitative/cluster_4/image_2.png}
        {Cluster 3: the ball is near the top of the field, close to the player's paddle.}
    \qualitativeClusterTwo
        {figures/clustering_agent_representation/pictures/pong/qualitative/cluster_2/image_1.png}
        {figures/clustering_agent_representation/pictures/pong/qualitative/cluster_2/image_2.png}
        {Cluster 4: same as cluster 1.}
    \qualitativeClusterTwo
        {figures/clustering_agent_representation/pictures/pong/qualitative/cluster_3/image_1.png}
        {figures/clustering_agent_representation/pictures/pong/qualitative/cluster_3/image_2.png}
        {Cluster 5: same as cluster 2.}
    \qualitativeClusterTwo
        {figures/clustering_agent_representation/pictures/pong/qualitative/cluster_5/image_1.png}
        {figures/clustering_agent_representation/pictures/pong/qualitative/cluster_5/image_2.png}
        {Cluster 6: same as cluster 3.}
    \caption{Clusters identified by \gls{car} in the \PongName environment (layer 4, $K=6$). Clusters 4--6 duplicate the semantics of clusters 1--3, illustrating over-segmentation at this value of $K$.}
    \label{fig:qualitative_pong}
\end{figure}

\vfill
